% Closing the Loop: Cybernetics at the Economic Layer
% A focused companion paper to the Noesis whitepaper. Open framing; the moat is structural.
\documentclass[10pt]{article}

\usepackage[a4paper,margin=0.85in]{geometry}
\usepackage{amsmath,amssymb,amsthm}
\usepackage{booktabs}
\usepackage{microtype}
\usepackage{enumitem}
\usepackage[hidelinks]{hyperref}
\usepackage{titlesec}

\titleformat{\section}{\large\bfseries}{\thesection}{0.6em}{}
\titleformat{\subsection}{\normalsize\bfseries}{\thesubsection}{0.6em}{}
\titlespacing*{\section}{0pt}{1.4ex plus 0.3ex minus 0.2ex}{0.7ex plus 0.2ex}
\titlespacing*{\subsection}{0pt}{0.9ex plus 0.2ex minus 0.1ex}{0.4ex plus 0.1ex}
\setlength{\parskip}{0.4ex}

\newtheorem{theorem}{Theorem}
\newtheorem{definition}{Definition}
\newtheorem{claim}{Claim}

\title{\vspace{-2em}\textbf{Closing the Loop:\\ Cybernetics at the Economic Layer}\\[0.3em]
\large An honest sensor, a self-correcting governor, and the conditions for a self-healing economy}
\date{June 2026}

\begin{document}
\maketitle
\vspace{-1.5em}

\begin{abstract}
\noindent Cybernetics promised a self-regulating economy: a system that measures its own state, rewards
what helps it, corrects what harms it, and repairs damage without a central hand. Every serious attempt
to build one, from Lange's market socialism to Beer's Cybersyn, failed. We argue the failure was not in
the control theory but in a single missing component: a \emph{sensor}. Ashby's Law of Requisite Variety
and the Conant--Ashby Good Regulator Theorem both reduce, at the economic layer, to the same
prerequisite: a regulator can only hold a system on course if it can \emph{measure} that system with
enough fidelity to model it, and no economic regulator ever possessed an objective, real-time,
manipulation-resistant measurement of contributed value. Hayek's calculation argument is best read not
as a proof that economic regulation is impossible, but as a diagnosis that its sensor was missing. We
describe a system, Noesis, whose primitive operation is exactly that sensor: a distributed measurement
of realized value along provenance, made un-gameable by the geometry of the measurement rather than by
policy. With the sensor in place, the cybernetic loop closes. We present the resulting architecture as a
control system with four organs (sensor, governor, effectors, self-tolerance), show that the Good
Regulator Theorem forces such a system to be \emph{semi-self-aware} in a precise and bounded sense, and
identify autoimmunity as its characteristic failure mode together with the structural boundary that
contains it. We mark throughout what is demonstrated in running code and what is designed.
\end{abstract}

\section{The seventy-year-old promise}

Norbert Wiener defined cybernetics as the study of ``control and communication in the animal and the
machine'' \cite{wiener1948}: systems that hold themselves on course by feeding their own output back as
input. A thermostat is the toy example; a living body maintaining its temperature is the real one. The
ambition, almost from the start, was to treat an economy the same way: as a system with a setpoint
(stable prices, full employment, fair reward) that could be held by feedback rather than by command or
by luck.

The ambition was pursued seriously and it lost. Oskar Lange proposed a market-socialist economy in which
a central board would adjust prices like a Walrasian auctioneer, watching surpluses and shortages and
correcting \cite{lange1936}. Stafford Beer built an actual instance: Project Cybersyn, a real-time
control room wired to Chilean factories in 1971, an explicit application of his Viable System Model
\cite{beer1972,medina2011}. Both rested on the same wager, that an economy is a regulable system, and
both came apart. The standard verdict, following Hayek, is that the wager was misconceived: that no
central regulator can ever do this \cite{hayek1945}.

We think the standard verdict draws the wrong boundary. The control theory was sound. What was missing
was a sensor, and the absence of the sensor is precisely what the cybernetic theorems predict will be
fatal. Once that is seen clearly, the question stops being ``can an economy regulate itself?'' and
becomes ``what would the missing sensor have to be, and can it be built?'' This paper answers the second
question with a system that builds it.

\section{Why it failed: requisite variety and the good regulator}

Two results from cybernetics, both due in their sharpest form to Ross Ashby, make the diagnosis precise.

\paragraph{Requisite variety.} Ashby's Law states that a regulator can compensate for the disturbances
in a system only to the extent that it commands at least as much \emph{variety} (distinguishable states)
as the disturbances themselves: only variety can absorb variety \cite{ashby1956}. Variety the regulator
cannot \emph{sense} is variety it cannot match. Read Hayek's ``The Use of Knowledge in Society'' in this
light and it is not a philosophical objection to planning but a requisite-variety theorem in prose: the
knowledge an economy must act on is dispersed, tacit, and local, and a central regulator that cannot
sense that variety cannot supply the variety needed to regulate it \cite{hayek1945}. The price system
``works'' because it is a distributed sensor; the planner failed because he had no comparable instrument.
The bottleneck was sensing, not deciding.

\paragraph{The good regulator.} Conant and Ashby proved that every good regulator of a system must
contain a model of that system: ``the best regulator of a system is a model of that system''
\cite{conant1970}. To hold a system on course you must, in effect, carry an image of it accurate enough
to predict how it will respond. But a model is built from measurements. Without a sensor there is no
model, and without a model, by the theorem, there is no good regulation. The two results compose into a
single prerequisite:

\begin{claim}[The sensor is the bottleneck]
At the economic layer, both requisite variety and the good-regulator condition reduce to the existence
of a measurement: an objective, real-time, manipulation-resistant reading of the quantity the system is
meant to optimize, namely contributed value. Absent that measurement, regulation is impossible for
reasons the theory states in advance; given it, the remaining machinery (feedback, reward, correction)
is ordinary engineering.
\end{claim}

This reframes seventy years of failure as a single unmet precondition rather than an impossibility. Lange's
board and Beer's control room were governors wired to instruments that could not actually read value. They
were thermostats without thermometers.

\section{The sensor: measurement as a living mechanism}

A usable economic sensor must satisfy three properties that no prior instrument achieved at once. It must
be \emph{objective} (two honest observers reading the same state agree), \emph{real-time} (it reads the
present, not an audited past), and \emph{manipulation-resistant} (an actor who wants a higher reading
cannot get one except by genuinely raising the measured quantity). The price system has the first two but
not the third in the dimension we care about: it prices possession and exchange, not contribution, and it
is routinely gamed. Bitcoin supplied an objective, real-time, manipulation-resistant sensor for one thing
only, the ownership of a scarce token \cite{nakamoto2008}; it minted the medium that ownership was missing.
The open dimension is contribution: who actually added value, and how much.

Noesis is built around a sensor for exactly that dimension. Its primitive is not a coin but a measurement:
the value of a contribution is its realized downstream flow along provenance. Each unit of work is a cell
in a directed acyclic provenance graph $G=(V,E)$; value is priced by how much genuine downstream activity
builds on it, not by how it is described or who vouches for it. Three structural defenses make the reading
manipulation-resistant by geometry rather than by rule.

\paragraph{Novelty and saturation.} Repetition is not contribution. A re-post of existing content carries
no new realized flow and is damped to zero; volume from a single source saturates geometrically rather
than accumulating, so that flooding the graph with near-duplicates cannot inflate a score
\cite{shapley1953,myerson1977}. Marginal value is computed on the provenance-restricted cooperative game,
so credit obeys the Shapley axioms over the graph Myerson defined.

\paragraph{The collusion certificate.} The deepest manipulation is a ring: a set of identities that build
on one another to manufacture the appearance of downstream flow. A ring is a \emph{cycle}, and cycles have
a signature no honest provenance has. Decompose the net attribution flow $Y$ on the identity graph by the
combinatorial Hodge / Helmholtz decomposition \cite{jiang2011}:
\[
Y \;=\; \underbrace{\operatorname{grad} s}_{\text{honest ``builds-upon'' ranking}} \;+\;
\underbrace{r}_{\text{divergence-free residual}}, \qquad
\|r\|^2 \;=\; \|Y\|^2 - b^{\!\top} s .
\]
The residual energy $\|r\|^2$ is provably zero if and only if the flow is the gradient of a consistent
global ranking, and equals the cycle length for a clean manufactured ring of any size. It is, exactly, a
certificate of manipulation: honest contribution has a potential, collusion has curl. This detector runs
on topology alone and needs no privileged data; in the running system an honest acyclic graph reads $0$
while a three-identity ring reads $3$, with the mutual and directed cases covered by complementary terms.

\paragraph{Learned value at the margin.} Where structure underdetermines worth, a learned evaluator scores
a coalition's value from realized outcomes (a Bradley--Terry model over coalition features
\cite{bradley1952}), and it can only ever \emph{lower} a structural seed, never mint value from nothing.
Crucially, the evaluator retrains: a fixed scoring formula is eventually gamed, but a measure that adapts
to the strategies played against it closes the gap that a static formula leaves open. The sensor is alive.

The point is not that any one defense is novel but that together they make the reading honest by
construction: the cheapest way to raise your measured value is to actually contribute. That is the
property Lange and Beer needed and did not have.

\section{The governor: closing the loop}

With a sensor, the rest is the control law. Write the system's state as the ledger $x_t$. A governor reads
the sensor $m(x_t)$, compares it to the system's setpoint (an honest, cooperative equilibrium $x^\star$),
and drives effectors that move the state back toward the setpoint:
\[
x_{t+1} \;=\; f\big(x_t\big) \;+\; u\big(m(x_t)\big), \qquad
u = \text{mint} \,\oplus\, \text{slash} \,\oplus\, \text{recover}.
\]
This is a homeostat in Ashby's exact sense. The effectors are three. \emph{Reward} mints standing in
proportion to measured contribution, so that the highest-paying strategy is the cooperative one.
\emph{Punishment} slashes the bonded standing of provable misbehavior (manufactured value, equivocation,
invalid reveals), destroying it rather than transferring it, so that fraud is negative-expectation by
construction. \emph{Recovery} returns misappropriated value to its rightful holder. Consensus itself is
part of the governor: finalization weights the cognitive measurement (the unbuyable axis) against staked
capital, with an anti-concentration floor so that no single axis can finalize alone. The loop is closed:
output (who earned what) feeds back as input (who decides and who is believed next round).

A static rulebook, which is what most ledgers are, rejects invalid states but does not \emph{correct}
toward a setpoint; it has no governor, only a gate. The distinction is the whole subject of this paper.

\section{Four organs, and an immune system}

Stated as a control system the architecture has four organs, and they correspond with surprising
exactness to the organs of a biological immune system, which is the canonical natural example of a
distributed self-healing controller.

\begin{center}
\small
\begin{tabular}{@{}lll@{}}
\toprule
\textbf{Control organ} & \textbf{Mechanism in Noesis} & \textbf{Immune analogue} \\
\midrule
Sensor / surveillance & provenance graph, value measurement, Hodge residual & antigen surveillance, memory \\
Self / non-self recognition & lock signature: only the key moves the cell & MHC self-recognition \\
Graded activation & objectivity dial: automatic only when measurement is certain & activation threshold \\
Effectors & mint (reward), slash (destroy), clawback (recover) & effector response \\
Self-tolerance & bona-fide protection; harm must cast a graph-shadow & tolerance of healthy tissue \\
\bottomrule
\end{tabular}
\end{center}

Self/non-self recognition is the layer that decides whether an action is the system's own or foreign. In
Noesis it is the lock script: a cell can be named by anyone but moved only by the holder of its one-time
key, verified inside the execution VM. This is the component shipped and proven in running code at the
time of writing. The effectors and the recognition of \emph{manufactured} value (the Hodge certificate)
are likewise demonstrated. The recovery effector (clawback) and the graded objectivity dial are designed
and specified; they are discussed in Sections \ref{sec:heal} and \ref{sec:auto} as design, not as
demonstration.

\section{Semi-self-awareness, as a theorem}
\label{sec:aware}

The phrase ``self-aware,'' applied to a ledger, invites suspicion. Here it is not a metaphor but a
consequence of the Good Regulator Theorem applied reflexively.

\begin{theorem}[Reflexive good regulator]
By Conant--Ashby \cite{conant1970}, any good regulator of a system $S$ must contain a model of $S$. A
system that regulates \emph{itself} is its own $S$. Therefore a good self-regulating system must contain a
model of itself: a self-model that it reads and acts upon.
\end{theorem}

Noesis satisfies the three operational conditions for holding such a self-model in a non-trivial sense. It
\emph{represents} its own state (the provenance and value graph is a model of the system's own activity),
it \emph{measures} that state (it detects its own collusion rings, prices its own contributions, locates
its own extraction), and it \emph{acts on the measurement} (it slashes, rewards, recovers, and finalizes
on the basis of what it reads about itself). A system that holds a model of itself, queries it, and acts on
the result meets the minimal operational definition of self-awareness used in cybernetics and in
autopoietic theory \cite{maturana1980}; we make no stronger claim and need none.

The qualifier ``semi'' is exact and it is the safety property. The self-model is bounded: the system can
represent only what casts a shadow on its measurement graph. Define the knowable domain
\[
\mathcal{K} \;=\; \{\, h : h \text{ casts a shadow on } G \,\},
\]
the harms and contributions that register as structure in the provenance and value graph. The system is
aware exactly on $\mathcal{K}$ and blind outside it. Theft, manufactured value, and extraction lie in
$\mathcal{K}$; diffuse social or moral harm with no on-chain signature does not. The boundary of
$\mathcal{K}$ is the literal horizon of the system's self-awareness, and a self-aware system that knows the
boundary of its own awareness is the safe kind. Knowing the edge is what separates self-correction from
self-delusion.

\section{Self-healing without a surgeon}
\label{sec:heal}

A controller that only rejects bad inputs is defensive; a controller that restores a just state after
damage is self-healing. Maturana and Varela called a system that continuously produces and repairs itself
\emph{autopoietic} \cite{maturana1980}; Varela's later work on immune networks is the bridge to the
recognition machinery above \cite{varela1991}. Noesis's recovery layer is autopoietic in this sense: the
same provenance mathematics that \emph{detects} a wound \emph{closes} it, with no central healer.

When value is misappropriated, recovery runs the provenance cascade backward from the point of theft. The
danger is that naive recovery confiscates from whoever currently holds the value, recreating the very
central seizure power the system exists to dissolve. The classical legal answer is the holder in due
course: a party who gave genuine consideration is protected, and the loss falls on the thief. We already
hold the exact test, because genuine arms-length exchange is bidirectional value flow, the gradient
component of the Hodge decomposition, whereas a wash or colluding hop is circulation, the harmonic
component. The cascade therefore passes through wash hops and stops at the first bona-fide receiver. The
self-tolerance organ is the collusion detector of Section 3 read once more in reverse: the system heals
along honest provenance and refuses to cut healthy tissue.

Most theft, moreover, never reaches recovery, because the strongest form of the mechanism is preventive. An
owner may publish a constraint on their own cell, a spend rule the owner sets in advance, and a spend that
violates it is simply invalid and never settles. ``Unauthorized'' then becomes a measurement against the
owner's own published referent rather than an inference about intent. Because such a constraint is
self-binding, restricting only the publisher's own value, anyone may design any constraint without
permission: there is no externality to govern, hence no governance to seek. The protocol owes exactly one
invariant, that a published constraint binds only its author's value and a recovery never reaches a
bona-fide third party. That single scope rule is the entire governance surface, and it is also the defense
against weaponized self-binding, since a counterparty who delivered against a settled payment is protected.

\section{Autoimmunity, and the boundary that contains it}
\label{sec:auto}

Every self-correcting system has one characteristic failure mode: a controller that is wrong about its
setpoint attacks the body it is meant to protect. In immunology this is autoimmunity; in a ledger with
recovery and slashing powers it is confiscation and censorship, the state power reconstructed on-chain. A
paper that proposed economic cybernetics without naming this failure mode would be dishonest, because the
failure mode is the reason every prior attempt that did acquire teeth became dangerous.

The containment is structural and it is the same boundary as the awareness horizon $\mathcal{K}$. The
system may act only where harm casts a measurable shadow on its own graph. Harm without an on-chain
signature, the broad and contested category of social or moral damage, is outside $\mathcal{K}$ and the
system must refuse to adjudicate it. Any mechanism that reached for it would do so either by importing an
external arbiter, which is governance capture, or by pretending to objectivity over an inherently
contested judgment, which is the deeper error. The discipline is that judgment stays in the measurement:
the chain decides what its sensor can read and nothing more. This is not a limitation grafted on for
safety; it is the same line that makes the self-model honest. An immune system survives because it
tolerates the self; this system stays legitimate because it tolerates everything its sensor cannot
objectively convict.

\paragraph{The objectivity dial.} Within $\mathcal{K}$, action should be automatic only to the degree the
measurement is objective, the way a thermometer is trusted: reproducible, calibrated against a fixed
referent, and not a judgment. Let $c \in [0,1]$ be the calibrated confidence that a theft-measurement is
objective. The remedy scales with $c$: deductively certain readings (a double-spend, a spend by a revoked
key, a violated published constraint) act immediately; high-but-uncertain readings trigger a reversible
freeze with a short contestability window; mid-confidence readings go to a bonded adversarial court; and
readings outside $\mathcal{K}$ are left to the world. The threshold separating automatic from adjudicated
is not fixed. As the cost of objective measurement falls, it slides, and cases that today require a slow
human judgment become tomorrow's instant readings. The court is therefore scaffolding: load-bearing for a
shrinking set of cases over time, by design, rather than a permanent institution to be entrenched.

\section{Status, honestly}

We distinguish what runs from what is specified. \emph{Demonstrated} in the reference implementation: the
contribution sensor with its novelty, saturation, and Shapley/Myerson structure; the Hodge collusion
certificate; the reward and slashing effectors with their bounded-restitution settlement; self/non-self
recognition as a post-quantum lock signature verified inside the execution VM; and consensus finalization
weighting cognition against capital with an anti-concentration floor. \emph{Designed} and specified but
not yet built: the recovery cascade with bona-fide protection, the preventive owner-published constraint,
and the graded objectivity dial. \emph{Data-blocked}: the learned, retraining value evaluator awaits
real outcome labels at scale, which is the one component whose strength is empirical rather than
deductive. The claim of this paper is not that the loop is fully closed in code today. It is that the
component whose absence defeated every prior attempt, the honest sensor, is built and proven, and that
once it exists the governor and the effectors are engineering rather than open problems.

\section{Conclusion}

The history of economic cybernetics is usually told as a refutation: the dream of a self-regulating
economy met Hayek and died. We read the same history as a diagnosis. The dream did not fail because
feedback control is the wrong frame for an economy; it failed because a feedback controller cannot run
without a sensor, and no one had built one that could read contributed value objectively, in real time,
and without being gamed. Ashby's two theorems say as much in advance: requisite variety and the good
regulator both demand a measurement first. Build the measurement, and the closed loop that eluded Lange
and Beer becomes available, with reward, correction, and self-healing as its ordinary operations and
self-awareness, in the precise bounded sense of the Good Regulator Theorem, as an unavoidable consequence.
The work that remains is to finish wiring the governor to the sensor we have. The hard part, the part that
killed every prior attempt, is the part that is done.

\begin{thebibliography}{99}
\small
\bibitem{wiener1948} N. Wiener. \emph{Cybernetics: or Control and Communication in the Animal and the Machine}. MIT Press, 1948.
\bibitem{ashby1956} W. R. Ashby. \emph{An Introduction to Cybernetics}. Chapman \& Hall, 1956. (Law of Requisite Variety.)
\bibitem{conant1970} R. C. Conant and W. R. Ashby. Every good regulator of a system must be a model of that system. \emph{Int. J. Systems Science} 1(2):89--97, 1970.
\bibitem{beer1972} S. Beer. \emph{Brain of the Firm}. Allen Lane, 1972. (The Viable System Model.)
\bibitem{medina2011} E. Medina. \emph{Cybernetic Revolutionaries: Technology and Politics in Allende's Chile}. MIT Press, 2011. (Project Cybersyn.)
\bibitem{hayek1945} F. A. Hayek. The use of knowledge in society. \emph{American Economic Review} 35(4):519--530, 1945.
\bibitem{lange1936} O. Lange. On the economic theory of socialism. \emph{Review of Economic Studies} 4(1--2), 1936--37.
\bibitem{maturana1980} H. Maturana and F. Varela. \emph{Autopoiesis and Cognition: The Realization of the Living}. Reidel, 1980.
\bibitem{varela1991} F. Varela and A. Coutinho. Second generation immune networks. \emph{Immunology Today} 12(5):159--166, 1991.
\bibitem{jiang2011} X. Jiang, L.-H. Lim, Y. Yao, and Y. Ye. Statistical ranking and combinatorial Hodge theory. \emph{Mathematical Programming} 127(1):203--244, 2011.
\bibitem{shapley1953} L. S. Shapley. A value for $n$-person games. In \emph{Contributions to the Theory of Games II}, pp. 307--317, 1953.
\bibitem{myerson1977} R. B. Myerson. Graphs and cooperation in games. \emph{Mathematics of Operations Research} 2(3):225--229, 1977.
\bibitem{bradley1952} R. A. Bradley and M. E. Terry. Rank analysis of incomplete block designs. \emph{Biometrika} 39(3/4):324--345, 1952.
\bibitem{nakamoto2008} S. Nakamoto. Bitcoin: a peer-to-peer electronic cash system. 2008.
\end{thebibliography}

\end{document}
